% bridge.tex
% ============================================================================
% NEW connective material -- the only wholly new prose in the combined
% document.  It joins Part I (the embedder, body of
% simplicial_complex_embedders_v3) to the injectivity proof that follows
% (body of 2026-Injectivity-Simplex-Method/main.tex).
%
% The dictionary table below maps the proof's notation onto the embedder's.
% (An earlier note here flagged an i/j axis "seam" between the two halves; it
% has been removed deliberately.  The i/j index difference is harmless --
% indices may be relabelled freely provided they agree on each side of an
% equation -- and the Part I text has been adjusted so nothing hinges on a
% physical reading of the row/column indices.)
% ============================================================================

\section{Injectivity: why local canonicalisation loses nothing}
\label{sec:bridge}

Section~\ref{sec:goals} asked for four properties.  Steps~A--F deliver
multiset-invariance, $C^0$ continuity and degree-1 homogeneity by construction,
and the discussion around Step~D made the continuity mechanism explicit (the
``greyed cells'': basis vectors that are ambiguous at a tie are precisely the
ones whose coefficient is zero there).  \emph{Injectivity} was the one
desideratum left on a heuristic footing.  This section discharges it---all but
one elementary step---and the rest of this section supplies the proof.

\subsection*{What injectivity reduces to}

The output (Step~F) is assembled from the multiset of pairs
$\bigl(\alpha_s,\ \canon{L_s}\bigr)$ together with the $k$ anchors $m_i$.  Grant
for the moment the Step~F cone mechanism: that $\sum_s\alpha_s\,\hash{\canon{L_s}}$
recovers, generically, the multiset of (coefficient, basis-vector) pairs that
built it---the ``generically disjoint $m$-cones in $\R^{2m+1}$'' argument of
Step~F, which we take here as elementary.  End-to-end injectivity then rests on
a single combinatorial question:

\begin{quote}
Do the \emph{locally} canonicalised basis vectors $\{\canon{L_s}\}$---each
$L_s$ column-sorted \emph{on its own} in Step~E---still determine the input
multiset $M$ up to the one global relabelling of its $n$ columns that
``multiset'' allows?
\end{quote}

The worry is real and invisible to numerical search.  Step~E sorts the columns
of each $L_s$ independently, throwing away which column of one $L_s$ corresponds
to which column of another.  A priori, two genuinely different multisets might
produce families $\{L_s\}$ that match \emph{matrix-by-matrix} after independent
sorting, yet are \emph{not} related by any single, common column permutation.
Such a collision would break injectivity, and---being a coincidence among
sorted integer matrices---would live on a measure-zero set that random probing
never hits.

\subsection*{The theorem, in the embedder's language}

The proof that fills out the rest of this section rules that out.  Restated in the
notation of Part~I:

\begin{tcolorbox}[keybox]
\textbf{Local canonicalisation is injective-safe.}
If, for two inputs, every locally canonicalised basis vector $\canon{L_s}$ on
one side equals one on the other (Step~E agreement), then a \emph{single}
column permutation $\sigma$ carries the whole family $\{L_s\}$ of one input onto
that of the other (\emph{global} agreement).  This is
Theorem~\ref{thm:main}; the operative direction is
\[
\text{local (per-matrix) agreement}\ \Longrightarrow\ \text{one global }\sigma.
\]
A single common $\sigma$ on the $n$ columns is exactly one relabelling of the
$n$ vectors---i.e.\ the two inputs are the \emph{same multiset}.  So Step~E
cannot merge distinct multisets.
\end{tcolorbox}

\noindent
The surprise is that the \emph{weaker}, each-matrix-separately notion of
agreement that Step~E actually tests already \emph{forces} the stronger,
all-at-once one.  That implication---local $\Rightarrow$ global---is precisely
the direction injectivity needs, and it is the hard direction of
Theorem~\ref{thm:main} (the easy converse, global $\Rightarrow$ local, is not
what we are after).

\subsection*{Dictionary}

The proof is written self-containedly in combinatorial language; every object
in it is one the embedder manufactures.  The correspondence is exact:

\begin{center}
\small
\setlength{\tabcolsep}{6pt}
\renewcommand{\arraystretch}{1.3}
\begin{tabular}{@{}p{6.1cm}p{5.0cm}p{3.6cm}@{}}
\toprule
\textbf{Proof object (the part below)} & \textbf{Embedder object (Part~I)} & \textbf{What it is} \\
\midrule
$\vec v_j\in s(n)$, leave-one-out order of $n$ symbols
  & descending sort order of row~$i$; left-out symbol $=$ that row's $\argmin$
  & per-row rank order \\
$coh(\vec v_j)_m$, indicator of prefix $P_m$
  & $V^{(i)}_{r}$, the top-$(r{+}1)$ set (output of Step~C, 1st BSD)
  & nested prefix indicator \\
$U(\vec v_1,\dots,\vec v_k)$, the $(n{-}1)k$ pieces
  & the $m=k(n{-}1)$ prefix basis elements pre-merge
  & the piece-set \\
ordering $p$ of $\vec U$
  & the gap-descending sort of Step~D
  & merge order \\
$C\vec U$, cumulative; $S^{\vec v}(Q)$, snapshot
  & $L_s=V_{(0)}+\dots+V_{(s)}$, the Eji\_LinComb count matrix (Step~D, 2nd BSD)
  & cumulative basis matrix \\
position $=\sum_r\max_c(\cdot)_{rc}$
  & the stored index $s{+}1$ of $L_s$
  & step count \\
filter $\vec f$ then $clean$
  & dropping the $\alpha_s=0$ terms
  & which snapshots survive \\
$canon$ / $localSort$ (Step~E analogue)
  & Step~E: sort each count matrix's columns on its own
  & local canonicalisation \\
$globalSort$
  & (not in the algorithm) one common $\sigma$ for all $L_s$
  & multiset-faithful sort \\
$\Sigma(Q)$
  & column relabellings aligning the prefixes a snapshot reveals
  & alignment set \\
\bottomrule
\end{tabular}
\end{center}

\subsection*{Why the proof carries a filter, two orderings and a ``clean''}

These three apparent complications in the statement of Theorem~\ref{thm:main}
are not decoration: each models exactly one feature of how the embedder behaves
at ties, and discharging them is what the proof spends its effort on.

\emph{The filter.}  The cone sum sees only nonzero-coefficient terms, so a pair
with $\alpha_s=0$ is not dropped by choice---it is \emph{invisible}.  Which $s$
survive is data-dependent (any pattern can occur, since
$\alpha_s=(s{+}1)(\delta_{(s)}-\delta_{(s+1)})$ vanishes wherever two merged
gaps coincide), so the safety statement must hold for \emph{every} filter
$\vec f$.  Moreover $\alpha_s=0$ holds exactly when $L_s$ sits inside a run of
equal gaps---exactly when $L_s$ is a \emph{tie-break-dependent} partial sum, the
ambiguous basis vector of the ``greyed cells''.  Hidden $\iff$ zero coefficient
$\iff$ ambiguous: the very vectors continuity \emph{forces} us to hide are the
ones whose value we could not have pinned down anyway, and the survivors are the
run-boundary sums, which are tie-break-invariant.  The theorem's job is to
prove this forced, data-dependent hiding is nonetheless lossless.

\emph{Two orderings $p,q$.}  Two nearby inputs straddling a tie break it in
opposite ways---different orderings of a run's pieces.  The proof allows
$p\neq q$ and shows local agreement \emph{forces them back into step}, so the
arbitrary tie-break costs nothing.

\emph{The $clean$.}  The cone sum is both order-blind and zero-blind, so hidden
positions leave no trace: the survivors arrive as a bare set with holes.  The
position-recovery lemma (Lemma~\ref{lem:recover}; row-maxima sum $=$ index)
lets each survivor re-announce where it sat, refilling the holes---which is why
$clean$ may be undone and the families compared as plain sets.

\begin{tcolorbox}[notebox]
\textbf{Scope.}  This closes the \emph{discrete} half of injectivity: the
identities of the surviving basis vectors, and that local canonicalisation
recovers them faithfully up to one global relabelling.  The remaining half---
that the Step~F cone sum recovers the (coefficient, basis-vector) pairs in the
first place---is the elementary generically-disjoint-cones argument flagged
above, and is not re-treated here.
\end{tcolorbox}

\subsection*{Reading on}

The rest of this section makes good on Theorem~\ref{thm:main}.  It
rebuilds the construction in self-contained combinatorial language
(Section~\ref{sec:defs}), isolates the precise question
(\S\ref{sec:problem}), and proves the local-equals-global equivalence from
first principles, the whole turning on a single idea---\emph{a snapshot
remembers how it was built}---read once shallowly (to neutralise $clean$) and
once fully (to neutralise $p\neq q$), with monotonicity then upgrading
``each snapshot is alignable'' to ``one $\sigma$ aligns them all''
(\S\ref{sec:solution}).  A reader who accepts the dictionary above may skim from
\S\ref{sec:problem} and rejoin at the worked examples that close the proof.
